\documentclass[a4j,10pt,onecolumn,oneside,final]{jreport}

%\usepackage{gthesis}
\usepackage{caption}
\usepackage{mthesis}

\usepackage{amsmath,amsthm,amssymb,ascmac}
\usepackage{nccmath}
\usepackage{fancybox}
\usepackage{slashbox}
%\usepackage[dviout]{color,graphics}
\usepackage{graphicx}
\usepackage{color}
\usepackage[a4paper,top=20truemm,bottom=20truemm]{geometry}
%\usepackage{layout}
\usepackage{psfrag}
%\usepackage{upgreek}
\usepackage{bm}
\usepackage{wrapfig}
\usepackage{ulem}
\usepackage{here}
%\usepackage[japanese]{babel}
\usepackage{rmath}

%\makeatletter
%\def\@cite#1{\textsuperscript{#1)}}
%\makeatother

%%%%%%%%%%%%%%%%%%%%%%%%%%%%%%%%%%%%%%%%%%%%%%%%%%%%%%%%%%%%%%%%
% put local tex-macros in this file %

%%%%%%%%%%%%%%%%%%%%%%%%%%%%%
%色設定
%%%%%%%%%%%%%%%%%%%%%%%%%%%%%
\definecolor{Black}{rgb}{0.0,0.0,0.0}
\definecolor{Red}{rgb}{0.9,0.0,0.1}
\definecolor{Blue}{rgb}{0.1,0.1,0.5}
\definecolor{Green}{rgb}{0.1,0.4,0.1}
\definecolor{Gray}{rgb}{0.75,0.85,0.9} % Fuji-iro
\definecolor{Shade}{rgb}{0.1,0.1,0.4}

%%%%%%%%%%%%%%%%%%%%%%%%%%%%%%%%%%%%%%%%%%%%%%%%%%%%%%%%%%%%%%%%
\年度{2020年度}
%\提出年月{平成29年7月}
\提出年月{2020年1月}

\題名{学位論文の\\体裁に関する賢雄}
%\梗概題名{学位論文の\\体裁に関する研究}
%\題名{楽に卒業する方法に関する基礎的研究\\
%			--- そもそもそれは可能か？ ---}
%\梗概題名{楽に卒業する方法に関する基礎的研究\\
%			--- そもそもそれは可能か？ ---}

\指導教員名A{情報 一郎}
\職名A{}

\指導教員名B{工学 次郎}
\職名B{}

%\所属学科{電気電子工学科}{}
\所属学院{情報理工学院}{知能情報コース}
\学籍番号{18M00000}
\氏名{東工大 太郎}

\概要{これはアブスト。
これはアブスト。
これはアブスト。
これはアブスト。
これはアブスト。
これはアブスト。
これはアブスト。
これはアブスト。
これはアブスト。
これはアブスト。
これはアブスト。
これはアブスト。
これはアブスト。
これはアブスト。
これはアブスト。
これはアブスト。
これはアブスト。
これはアブスト。
これはアブスト。
これはアブスト。
これはアブスト。
これはアブスト。
これはアブスト。
これはアブスト。
これはアブスト。
これはアブスト。
これはアブスト。
これはアブスト。
これはアブスト。
これはアブスト。
これはアブスト。
これはアブスト。
これはアブスト。
これはアブスト。
これはアブスト。
これはアブスト。
これはアブスト。
これはアブスト。
これはアブスト。
これはアブスト。
これはアブスト。
これはアブスト。
これはアブスト。
これはアブスト。
これはアブスト。
これはアブスト。
これはアブスト。
これはアブスト。
これはアブスト。
これはアブスト。
これはアブスト。
これはアブスト。
これはアブスト。
これはアブスト。
これはアブスト。
これはアブスト。
これはアブスト。
これはアブスト。
これはアブスト。
これはアブスト。
これはアブスト。
これはアブスト。
これはアブスト。
これはアブスト。
これはアブスト。
これはアブスト。
これはアブスト。
これはアブスト。
これはアブスト。
これはアブスト。
これはアブスト。
これはアブスト。
これはアブスト。
これはアブスト。
これはアブスト。
これはアブスト。
これはアブスト。
これはアブスト。
これはアブスト。
これはアブスト。
これはアブスト。
これはアブスト。
これはアブスト。
これはアブスト。
これはアブスト。
これはアブスト。
これはアブスト。
これはアブスト。
これはアブスト。
これはアブスト。
これはアブスト。
これはアブスト。
これはアブスト。
これはアブスト。
これはアブスト。
これはアブスト。
これはアブスト。
これはアブスト。
これはアブスト。
これはアブスト。
これはアブスト。
これはアブスト。
これはアブスト。
これはアブスト。
これはアブスト。
これはアブスト。
これはアブスト。
これはアブスト。
これはアブスト。
これはアブスト。
これはアブスト。
これはアブスト。
これはアブスト。
これはアブスト。
これはアブスト。
これはアブスト。
これはアブスト。
これはアブスト。
これはアブスト。
これはアブスト。
これはアブスト。
これはアブスト。
これはアブスト。
これはアブスト。
}


\begin{document}
% ----------------------------------------------------------------------
% 表紙
\maketitle
% ----------------------------------------------------------------------
% 目次
\setcounter{page}{1}
%\renewcommand{\thepage}{\roman{page}}
\renewcommand{\thepage}{\arabic{page}}  %% いる？
\setcounter{tocdepth}{1}
\tableofcontents
\clearpage
% ----------------------------------------------------------------------
% 内容
%\setcounter{page}{1}
\renewcommand{\thepage}{\arabic{page}}
\chapter{Introduction}

I introduce a sample of master's thesis \cite{taro}.


\section{Sample}

The followings are samples.
I assume an $n$-dimensional hypersphere shown in Eq. \eqref{eq:eq_1_1}.
In case of $n=3$, the shape is like in Figure \ref{fig:ball 1}.
\begin{align}
	r^2 = \sum_{k=1}^n x_k^2
	\label{eq:eq_1_1}
\end{align}
\begin{figure}[h]
	\centering
	\includegraphics[bb=0 0 79 65]{./figures/ball.png}
	\caption{A sphere where $n=3$}
	\label{fig:ball 1}
\end{figure}

According to Table \ref{table:elem 1}, it has four elements (a,b,c,d).
\begin{table}[h]
	\centering
	\begin{tabular}{|c|c|}
		\hline
		a & b \\ \hline
		c & d \\ \hline
	\end{tabular}
	\caption{Elements}
	\label{table:elem 1}
\end{table}


%\input{body-a}
%\input{body-b}
%\input{body-c}
\chapter{結論}

結論は、網羅的にかつ簡潔に\cite{hanako}。

\chapter*{Acknowledgment}
\addchapter{Acknowledgment}

Thank you.


\bibliographystyle{mthesis}
\bibliography{bibliography}
\addchapter{References}




% for those who don't use bibtex 
% \begin{thebibliography}{99}
% \addchapter{References}
%  	\bibitem{taro} Taro Tokodai, ``How to Write a Good Thesis'',
%  		{\it Journal of XYZ},
%  		Vol.3, No.4, pp.15--34,
%  		2015
% \end{thebibliography}

%\input{appendix1}
%\input{appendix2}
% ----------------------------------------------------------------------
\end{document}
